
% ##############################################################
% ##############################################################
% ##############################################################ermitteln
\chapter{Einleitung}


Dieses Buch stellt Ihnen alle wesentlichen Neuerungen aus den LTS-Releases
(Long Term Support)\index{LTS-Release}\index{Long Term Support}\index{LTS}
Java 17 LTS, Java 21 LTS und dem brandaktuellen Java 25 LTS vor.


Alle Leser, die privat oder beruflich noch auf die älteren Releases
Java 8 LTS oder Java 11 LTS setzen, finden in Anhang \ref{cha:Wesentliches aus Java 8 LTS, 9, 10 und 11 LTS}
einen kompakten Schnelleinstieg
zu den heraus\-ragendsten bzw. bedeutendsten Änderungen aus diesen  Java-Versionen,
um sich für die Erweiterungen in modernem Java inhaltlich und syntaktisch vorzubereiten.


Bevor wir uns mit den in diesem Buch behandelten Themen beschäftigen,
möchte ich die Veränderungen der letzten Jahre in der Releasepolitik von Oracle
darlegen, weil sich hier einiges getan hat.




% #########################################################
% #########################################################
% #########################################################
\section{Releasepolitik}
\index{Releasepolitik}


Vor Java 9 gab es keinen festen Takt oder Releasezyklus.\index{Releasezyklus}
Stattdessen war die Veröffentlichung einer neuen Java-Version an die Fertigstellung
wesentlicher Funktionen geknüpft. Viele erinnern sich bestimmt noch gut daran:
 In der Vergangenheit wurden Java-Releases aufgrund unfertiger Features häufiger verschoben.


Beispielsweise wurde Java 6 im Dezember 2006 veröffentlicht, während sein Nachfolger Java 7
erst im Juli 2011, also fast fünf Jahre später, erschien. Java 8 LTS, das im März 2014 veröffentlicht
wurde, verzögerte sich mehrmals, hauptsächlich aufgrund unvollendeter wichtiger Funktionalitäten.
Selbst Java 9 wurde mehr als einmal gegenüber dem ursprünglichen Zeitplan verschoben --
vor allem wegen der Komplexität der Modularisierung.


Diese langen und unvorhersehbaren Veröffentlichungsintervalle erschwerten sowohl den Anwendern als auch den Tool-Anbietern die Planung. Um dem entgegenzuwirken, hat Oracle nach der Veröffentlichung von Java 9 im September 2017 auf einen halbjährlichen Release-Zyklus umgestellt. Das erlaubt es,  die jeweils bis zu diesem Zeitpunkt fertig implementierten Funktionalitäten zu veröffentlichen.
Allerdings sind logischerweise nicht immer schon alle Features auch so weit, um final
veröffentlicht zu werden. Daher gibt es seit Java 9 zwei Besonderheiten:
die  Preview-Features und  die Incubator-Features.




% #########################################################
\subsubsection{Preview-Features}\index{Preview-Feature}
\label{subsubsec:Preview-Features}


Sogenannte Preview-Features sind zwar vollständig spezifiziert und implementiert.
Sie werden aber zunächst als Vorschau zum Sammeln von Erfahrungen und Feedback ins
JDK integriert. Basierend auf den Rückmeldungen kann sich die Implementierung,
etwa verwendete Typen, Methodensignaturen usw., allerdings noch ändern.
Insgesamt möchte man der Funktionalität durch die gesammelten Erfahrungen und Rückmeldungen
dann in Folgereleases einen Feinschliff  geben. In Ausnahmefällen können die gewonnenen Erkenntnisse auch zu
 einem Abbruch der Weiterentwicklung oder einer Neuausrichtung führen.


Preview-Features  sind standardmäßig nicht zugänglich
und müssen explizit über die Angabe von \code{-}\code{-enable-preview}
sowohl beim Kompilieren als auch beim Ausführen  aktiviert werden, um
sie in eigenen Programmen verwenden zu können.
Beim Kompilieren  muss (bis Java 22) zudem eine Source-Version
 mit  \code{-}\code{-source} angegeben werden --
seit Java 23 ist dies optional.




% #########################################################
\subsubsection{Incubator-Features}\index{Incubator-Feature}


Ergänzend zu den Preview-Features  gibt es  sogenannte Incubator-Features,
die in Form korrespondierender Module implementiert und bereitgestellt werden.
Auch hierbei geht es um das Sammeln von Erfahrungen und Feedback, allerdings
auf Basis einer  vorläufigen Umsetzung.
Es kann bei Incubator-Features durchaus sein, dass sich Dinge noch grundlegend ändern
oder Funktionalitäten sogar später vollständig entfernt
bzw. gar nicht in ein finales JDK aufgenommen werden. Weil Incubator-Features bzw.
die bereitstellenden Module nicht offizieller Bestandteil des JDKs sind, müssen diese speziell
eingebunden werden. Dazu dient der JVM-Aufrufparameter \code{-}\code{-add-modules}.
Dieser ist sowohl beim Kompilieren als auch beim Ausführen anzugeben.\footnote{Interessanterweise
gilt dies für das Kompilieren nicht innerhalb von Eclipse. Dort muss
lediglich die Run Configuration passend parametriert werden.}






% #########################################################
\subsubsection{Releasekadenz im Wandel und die Auswirkungen}


Kommen wir auf die Releasekadenz zurück.\index{Releasekadenz}\index{Releasekadenz!im Wandel}
Bis einschließlich Java 9 wurden neue Java-Versionen immer Feature-basiert
veröffentlicht.\index{Releasestrategie!Feature-basierte}
Das hatte in der Vergangenheit oftmals und mitunter auch beträchtliche  Verschiebungen des
geplanten Releasetermins zur Folge, nämlich immer dann, wenn  für die Version wesentliche Features noch
nicht fertiggestellt waren. Insbesondere deshalb verzögerten sich die Veröffentlichungen von
Java~8 LTS und Java~9 um mehrere Monate bzw. sogar über ein Jahr.
 Schlimmer noch: Es war damals absolut nicht abzusehen
oder vorab zu planen, wann eine neue Java-Version tatsächlich veröffentlicht werden würde.

Die Umstellung auf eine zeitbasierte Releasestrategie\index{Releasestrategie!zeitbasierte}
wirkt derartigen Verzögerungen entgegen und erhöht die Planbarkeit, weil jedes halbe Jahr
eine neue Java-Version veröffentlicht wird, die all jene Features enthält, die bis zu
dem Termin fertiggestellt sind.
Zudem sollte alle drei Jahre eine LTS-Version (Long Term Support) erscheinen. Eine
solche ist in etwa vergleichbar mit den früheren Major-Versionen. Das war zwischen Java 11 LTS
und 17 LTS  auch der Fall. Mit Letzterem wurde dann eine Verkürzung des
LTS-Zyklus auf 2 Jahre beschlossen, wodurch auch diejenigen Nutzer schneller von den
Neuerungen profitieren, die lediglich von LTS zu LTS springen wollen oder können.

Lassen Sie mich noch auf Folgendes hinweisen:
Zwar kann die schnelle, halbjähr\-liche Releasefolge eine größere Herausforderung für Toolhersteller sein,
für uns als Entwickler ist sie aber oftmals positiv, weil wir weniger lang auf neue
Features warten müssen.  Das konnte früher recht zermürbend sein.


Allerdings gilt das Positive vor allem für eigene Hobbyprojekte, weil man dort
mit den Neuerungen experimentieren kann und weniger durch Restriktionen
eingeschränkt ist.
Im professionellen Einsatz wird man eher auf Kontinuität und die Verfügbarkeit
von Security Updates setzen, weshalb in diesem Kontext vermutlich nur
LTS-Versionen in Betracht kommen, um Migrationsaufwände kalkulierbar und
 besser planbar zu halten.




% #######################################################################
\begin{praxistipp}{Am Rande ...}
\begin{center}
\begin{minipage}{\textwidth-0.6cm}
Für jeden Entwickler (Amateur oder Profi) hat der halbjährliche Releasezyklus\index{Releasezyklus!halbjährlicher}
einen weiteren Vorteil:
Man kann sich häppchenweise mit neuen Features beschäftigen, wenn sie erscheinen, und
muss nicht alle zwei, drei oder mehr Jahre einen Riesenbatzen Neuerungen am Stück durcharbeiten,
wobei man  durch den schieren Umfang an Änderungen gegebenenfalls einige relevante Dinge übersehen
kann oder  aus Zeitgründen auf \frqq das les ich mir später durch\flqq{} verschieben muss.
\end{minipage}
\end{center}
\end{praxistipp}




% #########################################################
\subsubsection{Lizenzpolitik im Wandel und die Auswirkungen}\index{Lizenzpolitik}\index{Lizenzpolitik!im Wandel}


Nicht nur beim Releasezyklus, sondern insbesondere auch bei der Lizenzpolitik
gab es in den letzten Jahren ein paar Überraschungen.


All diejenigen, die Java 11 LTS herunterladen wollten, wurden durch einen
markanten Hinweis auf die Änderungen in der Lizenzpolitik hingewiesen.


% 1?708×845, 150dpi => *.72
\begin{figure}[htb]
\centering
\includegraphics[bb=0 0 1229.76 608.4,width=\textwidth]{Abbildungen/NewLicense.png}
%\includegraphics[bb=0 0 1229.76 608.4,height=5.5cm]{Abbildungen/NewLicense.png}
\caption{Hinweis auf neue Lizenzbedingungen, exemplarisch für Java 11 LTS}
\end{figure}


\noindent
Dort war zu entnehmen, dass sich die Lizenzierungsbedingungen ändern.
Das hat Einfluss, wenn Sie Ihre Software  kommerziell  vertreiben oder dies planen zu tun:
\important{Das Oracle JDK war vor Java 8 LTS selbst in Produktions"-sys"-temen
immer kostenfrei verwendbar, wurde dann aber mit Java 11 LTS leider kostenpflichtig}.\index{OpenJDK}\index{Oracle JDK}\index{Lizenzpolitik}
Als Alternative existiert etwa das OpenJDK (\url{https://openjdk.org/}).
Während der Entwicklung kann das Oracle JDK allerdings weiterhin kostenfrei genutzt werden.


Dieser Schritt, die Benutzer zur Lizenzierung zu nötigen, hat bei der Entwickler\-gemeinde
 für Verstimmung gesorgt. Und es kam noch schlimmer: Kurze Zeit nach der Einschränkung für
Java 11 LTS wurden rückwirkend auch die Lizenzierungsbedingungen für Java 8 LTS
auf kostenpflichtig geändert. Viele Firmen haben daher erst einmal
keine weiteren Updates auf neuere Versionen vorgenommen.


Insgesamt kam es zu einer massiven Gegenreaktion. Diese wurde erleichtert,
weil das OpenJDK seit Java 11 LTS auf den gleichen Sourcen wie das Oracle JDK beruht
und man das OpenJDK kostenfrei beziehen kann. Basierend darauf haben sich diverse Firmen und
Konsortien daran gemacht, freie Versionen von Java bereitzustellen.
Hier seien Azul Zulu, Adopt OpenJDK (mittlerweile Adoptium Temurin\index{Adoptium Temurin}) oder
Amazon Corretto\index{Amazon Corretto} stellvertretend für viele weitere genannt.


\important{Bei Oracle hatte man mit der Veröffentlichung von Java 17 LTS ein Einsehen und das Oracle JDK
steht nun auch für kommerzielle Zwecke wieder kostenfrei zur Ver"-fügung}. Die zugehörige Lizenz
wird übrigens als No-Fee Terms and Conditions (NFTC)\index{NFTC}\index{No-Fee Terms and Conditions} bezeichnet.


Ein wenig unverständlich ist allerdings, dass Java 11 LTS weiterhin einer kostenpflichtigen Lizenz unterliegt.
Damit wird es noch unattraktiver und wohl kaum mehr als Zwischenschritt auf dem Weg zu Java 17 LTS,
 Java 21 LTS oder Java 25 LTS eingesetzt werden -- es sei denn, man nimmt ein anderes JDK als das von Oracle.




% ##############################################################
% ##############################################################
% ##############################################################
\section{Inhaltsübersicht}


Dieses Buch behandelt diverse Erweiterungen,
die in Java 17 LTS, Java 21 LTS sowie dem brandaktuellen Java 25 LTS gebündelt sind.
Diese LTS-Versionen werden jeweils in eigenen Teilen des Buchs beschrieben, sodass Sie
sehr schnell die wesentlichen Neuerungen der von Ihnen genutzten oder präferierten
Java-Version kompakt und fundiert im Überblick haben.




% ##############################################################
% ##############################################################
\subsection{Teil I: Neuerungen in Java 12 bis 17 LTS}


In diesem Teil sind die Neuerungen ein wenig kompakter und lediglich in einer
Auswahl beschrieben, da Java 17 LTS mittlerweile weitverbreitet sein sollte.
Für einen fundierten Einstieg und eine ausführlichere Darstellung der
relevanten Neuerungen verweise ich Sie auf mein
Buch \frqq Java 21 LTS bis 24 -- Coole neue Java-Features\flqq{} \cite{Java-21-LTS-24:2025}.


Dieses Kapitel fokussiert sich auf verschiedene Änderungen an der Syntax von Java, etwa
 \frqq Text Blocks\flqq{}, die mehrzeilige Strings erlauben.
Auch bezüglich \code{switch} hat sich einiges getan.
Die neue Syntax erleichtert die Angabe und Auswertung von Bedingungen und
sorgt für mehr Klarheit und Verständlichkeit.
Zusätzlich bietet modernes Java mit Records ein ganz besonders
 interessantes Feature, womit eine extrem kompakte Schreibweise zum Deklarieren
spezieller Klassen mit unveränderlichen Daten bereitgestellt wird.
Darüber hinaus gibt es einige für die Praxis weniger relevante Features, die aber teilweise
die Grundlage für weitere Neuerungen in zukünftigen Java-Versionen bilden:
 Im Kontext von \code{instanceof} ist es möglich,
künstliche Hilfsvariablen und unschöne Casts zu vermeiden.


Kommen wir zu den APIs: Hier finden wir  Ergänzungen in der Klasse \code{String}.
Zudem enthält das Stream-API einige nennenswerten Neuerungen, unter anderem den  Teeing-Kollektor
 sowie die Methoden \code{toList()} und  \code{mapMulti()}.
Des Weiteren wurde als JVM-Neuerung eine hilfreiche
Verbesserung zur Fehleranalyse bei \code{NullPointerException}s integriert.






% ##############################################################
% ##############################################################
\subsection{Teil II: Neuerungen in Java 18 bis 21 LTS}


Bezüglich der Syntax finden wir in den Java-Versionen 18 bis 21 LTS
diverse  wesentliche Neuerungen -- teilweise
allerdings erst als Preview-Feature. Mit String Templates ließ sich in Java 21 LTS und
Java 22 die Konkatenation von  variablen und fixen Textbestandteilen zu einem Ergebnisstring vereinfachen.
 Dieses Preview-Feature wurde mit Java 23 wieder aus dem JDK entfernt. Daher gehe
 ich nicht weiter darauf ein.

Record Patterns erweitern das Pattern Matching für \code{instanceof} und
ermöglichen es, Records auf ein\-fache Weise  in ihre Bestandteile zu zerlegen und darauf
zuzugreifen. Das Pattern Matching wurde ursprünglich im Kontext von
 \code{instanceof}\ in die Sprache integriert. Es funktioniert in modernem Java
 auch in \code{switch} und wurde funktionell aufgewertet.
Dabei ist insbesondere die Kombination mit den Record Patterns erwähnenswert.
 Schließlich wird es durch Unnamed Patterns and Variables möglich,
verschiedene Elemente in einem Record Pattern oder eine Variable durch
einen einzelnen Unterstrich (\frqq\code{\_}\flqq{}) zu ersetzen, um diese als ungenutzt und unbrauchbar zu markieren.


Java 21 LTS enthält einige API-Neuerungen, etwa kleinere
Erweiterungen rund um Reflection mit Method Handles und die Adressauflösung von
Internetadressen. Allerdings gibt es mit Sequenced Collections und
Virtual Threads auch zwei Highlights. Sequenced Collections erlauben es,
verschiedene Arten von Collections gleichartig von vorne und von hinten zu verarbeiten.
Mit Virtual Threads erzielt man eine enorme Vereinfachung beim Multithreading und es lassen sich
 jetzt mehrere Tausend, ja sogar Hunderttausende oder mehr
Threads erzeugen. Dadurch kann man dem Thread-per-Request-Ansatz folgen.


Ergänzend schauen wir uns einige Änderungen, Erweiterungen und Neuerungen in der JVM an,
die kumulativ in Java 18 bis 21 LTS enthalten sind. Das sind unter anderem der
 Simple Web Server sowie eine Verbesserung bei der Angabe von
  Code Snippets in JavaDoc. Abschließend wird das Feature
   Unnamed Classes and Instance Main Methods behandelt, das insbesondere die Art und Weise
   des Applikationsstarts verändert und zudem Vereinfachungen beim Einstieg in
   Java bringt, indem man auf Klassen und einigen Boilerplate-Code bei
   kleineren Java-Programmen verzichten kann.



% ##############################################################
% ##############################################################
\subsection{Teil III: Neuerungen in Java 22 bis 25 LTS}


Bezüglich der Syntax gibt es in den Java-Versionen 22 bis 25 LTS diverse  wesentliche Neuerungen -- teilweise
sind diese allerdings erst als Preview-Feature realisiert.


Um Variablen oder Teile innerhalb von Record Patterns durch einen Unterstrich (\frqq\code{\_}\flqq{})
als unbenutzt und unbrauchbar zu markieren, dient das finale Feature namens \frqq Unnamed Variables \& Patterns\flqq{}.
Als weiteres finales Feature ist das Kommentieren von Sourcecode mit Markdown integriert.
Damit kann man eine besser lesbare und knappere Dokumentation als mit den bisherigen JavaDoc-HTML-Schnipseln definieren. Ebenfalls finale Features sind die Module Import Declarations und die Flexible Constructor Bodies:
Die Imports am Anfang einer Java-Datei lassen sich durch Modul-Imports aufgeräumter  und übersichtlicher gestalten.
Flexible Constructor Bodies erleichtern Initialisierungsaktionen.
Noch als Preview-Feature findet sich die Neuerung Primitive Types in Patterns. Damit
lassen sich nun auch Pattern für primitive Typen in \code{instanceof} und \code{switch}
angeben.


Bezüglich API-Erweiterungen bringt Java 25 LTS beispielsweise
das Foreign Function \& Memory (FFM)  API\index{Foreign Function \& Memory APIs}. Dieses
stellt ein API zum Zugriff auf Funktionalitäten und Speicherbereiche, die außerhalb der JVM liegen, bereit.
Ebenso sind die Stream Gatherer  eine sehr nützliche Ergänzung. Diese erweitern das Stream-API um
die Möglichkeit, benutzerdefinierte Intermediate Operations bereitstellen zu können.
Praktischerweise sind bereits ein paar Stream Gatherer  vordefiniert.
 Als Highlight soll die Structured Concurrency nicht unerwähnt bleiben.
 Diese stellt eine massive Vereinfachung im Kontext von Multi\-threading dar und  erlaubt es,
 Aufgaben in Teil\-aufgaben aufzuspalten, diese  parallel zu verarbeiten und deren Ergebnisse  zusammenzuführen.
 Dabei existieren verschiedene Strategien, wie dabei vorgegangen werden soll, etwa für die Fälle, dass alle oder nur
 ein Ergebnis relevant ist.


Als JVM-Verbesserungen wurde unter anderem Launch Multi-File Source-Code Programs umgesetzt.
Das bietet die Möglichkeit, mehrere einzelne Java-Dateien direkt ohne explizites vorheriges Kompilieren ausführen.
Ergänzend erleichtern die Compact Source Files and Instance Main Methods den Einstieg in  Java.




% #########################################################
% #########################################################
\section{Grundgerüst des Beispielprojekts}\index{Beispielprojekt!mitgeliefertes}\index{Beispielprojekt!Grundgerüst}


Das mitgelieferte Projekt orientiert sich in seinem Aufbau an
dem  des Buchs und bietet für jedes Kapitel
jeweils ein eigenes Package, z.\,B. \code{ch04\_syntax\_java\_18\_21} oder
\code{ch05\_api\_java\_18\_21}. Dabei weiche ich ausnahmsweise von
der Namenskonvention für Packages ab, weil ich die Unterstriche
 in diesem Fall für eine besser lesbare Notation halte.

Das Projekt folgt dem Maven-Standardverzeichnisaufbau mit den
Sourcen unter \code{src/main/java} und den Tests unter \code{src/test/java}.\footnote{Da in diesem Buch vor allem
Beispiele zur Demonstration der Neuerungen enthält, verzichte ich auf die in der Praxis empfehlenswerten Unit Tests.
Daher ist das Testverzeichnis leer.}


%725?×?1219
\begin{figure}[htb]
\centering
\includegraphics[bb=0 0 522 877.68,height=11.5cm]{Abbildungen/J25b/Project-Java-25-LTS-src-main-java.png}
\end{figure}














\vfill
\pagebreak
% #########################################################
% #########################################################
% #########################################################
\section{Anmerkung zum Programmierstil}\index{Programmierstil}


In diesem Abschnitt möchte ich  vorab etwas zum Programmierstil
sagen, weil   bei Ihnen vielleicht ab und an einmal die Frage aufkommen mag,
ob man gewisse Dinge nicht kompakter gestalten könnte oder sollte.




%#######################################################################
%#######################################################################
\subsection{Gedanken zur Sourcecode-Kompaktheit}


In der Regel sind  mir  beim Programmieren und insbesondere für die
Implementierungen in diesem Buch vor allem eine leichte Nachvollziehbarkeit sowie
 eine übersicht\-liche Strukturierung und damit  später eine vereinfachte
  Veränderbarkeit wichtig. Deshalb sind die gezeigten Implementierungen
 möglichst verständlich programmiert.
 Dadurch ist vielleicht  nicht jedes Konstrukt maximal kompakt, dafür aber in
 der Regel  gut nachvollziehbar. Diesem Aspekt möchte ich in diesem
 Buch den Vorrang geben. Auch in der Praxis  kann man damit oftmals besser
 leben als mit einer schlechten Wartbarkeit, dafür aber einer kompakteren
 Realisierung.




% #########################################################
\subsubsection{Beispiel}


Zur Verdeutlichung betrachten wir eine lesbare, gut verständliche,
rekursive Implementierung zum Umdrehen des Inhalts eines Strings, die zudem sehr schön die
beiden wichtigen Elemente des rekursiven Abbruchs und Abstiegs zeigt:


\lstset{mathescape=true}
\vspace{-0.85mm}
\begin{lstlisting}
public String reverseString(final String input)
{
    // $\mbox{\bfseries rekursiver Abbruch}$
    if (input.length() <= 1)
        return input;


    final char firstChar = input.charAt(0);
    final String remaining = input.substring(1);

    // $\mbox{\bfseries rekursiver Abstieg}$
    return $\mbox{\bfseries reverseString(remaining)}$ + firstChar;
}
\end{lstlisting}
\vspace{-0.85mm}
\lstset{mathescape=false}


\noindent
Die folgende deutlich kompaktere Variante bietet diese
Vorteile nicht:


\lstset{mathescape=true}
\vspace{-0.85mm}
\begin{lstlisting}
public String reverseStringShort(final String input)
{
    return input.length() <= 1 ? input :
           reverseStringShort(input.substring(1)) + input.charAt(0);
}
\end{lstlisting}
\vspace{-0.85mm}
\lstset{mathescape=false}


Überlegen Sie  kurz, in welcher der beiden Methoden Sie
sich sicher fühlen, eine nachträgliche Änderung vorzunehmen.
Und wie sieht es aus, wenn Sie noch Unit Tests
ergänzen wollen: Wie finden Sie passende Wertebelegungen und Prüfungen?


Außerdem sollte man bedenken, dass die
obere Variante entweder bereits während der
Kompilierung (Umwandlung in den Bytecode) oder
später während der Ausführung und Optimierung
automatisch in etwas Ähnliches wie die untere Variante konvertiert wird --
eben mit dem Vorteil der besseren Lesbarkeit beim Programmieren.






%\vfill
%\pagebreak
%#######################################################################
%#######################################################################
\subsection{Gedanken zu \texttt{final} und \texttt{var}}


Normalerweise bevorzuge ich, Variablen als \code{final} zu markieren, um diese
als unveränderlich zu kennzeichnen. In diesem Buch verzichte ich mitunter darauf.
Ein Grund ist, dass die interaktive Kommandozeilenapplikation JShell (vgl. Abschnitt \ref{sec:Java + REPL})
das Schlüsselwort \code{final} nicht überall unterstützt, aber für Parameter und lokale Variablen
innerhalb von Methoden.


\subsubsection{Local Variable Type Inference -- \texttt{var}}\index{Local Variable Type Inference}\indexCode{var}
Seit Java 10 existiert die sogenannte \importantAndIndex{Local Variable Type Inference},
besser bekannt als \code{var} (vgl. Abschnitt \ref{sec:Syntaxerweiterung var}).
Diese erlaubt es, auf die explizite Typangabe
auf der linken Seite einer Variablendefinition zu verzichten, sofern  sich der
konkrete Typ für eine lokale Variable anhand der Definition auf der
rechten Seite der Zuweisung vom Compiler ermitteln lässt:


\lstset{morekeywords={var}}
\lstset{mathescape=true}
\vspace{-0.85mm}
\begin{lstlisting}
var name = "Peter";                // $\mbox{\bfseries var => String}$
var chars = name.toCharArray();    // $\mbox{\bfseries var => char[]}$
var mike = new Person("Mike", 47); // $\mbox{\bfseries var => Person}$
var hash = mike.hashCode();        // $\mbox{\bfseries var => int}$
\end{lstlisting}
\vspace{-0.85mm}
\lstset{mathescape=false}


\noindent
Insbesondere im Zusammenhang mit generischen Containern spielt die
Local Variable Type Inference ihre Vorteile aus:


\lstset{mathescape=true}
\vspace{-0.85mm}
\begin{lstlisting}
// $\mbox{\bfseries var => ArrayList<String> }$
var names = new ArrayList<String>();
names.add("Tim");
names.add("Tom");
names.add("Jerry");


// $\mbox{\bfseries var => Map<String, Long> }$
var personAgeMapping = Map.of("Tim", 47L, "Tom", 12L,
                              "Michael", 47L, "Max", 25L);
\end{lstlisting}
\vspace{-0.85mm}
\lstset{mathescape=false}


\subsubsection{Konvention: \texttt{var}, falls lesbarer}
Sofern die Verständlichkeit nicht darunter  leidet, werde ich \code{var}
an geeigneter Stelle verwenden, um den Sourcecode kürzer und
klarer zu halten. Ist jedoch eine Typangabe für das Nachvollziehen
von größerer Wichtigkeit, bevorzuge ich den konkreten Typ und
 vermeide \code{var} -- die Grenzen sind aber fließend.


\subsubsection{Konvention: \texttt{final} oder \texttt{var}}
Eine weitere  Anmerkung noch: Zwar kann man \code{final} und \code{var}
kombinieren, ich finde dies jedoch stilistisch nicht schön, weil dann die
Kürze von \code{var} verloren geht. Daher verwende ich
entweder das eine oder das andere.








%#######################################################################
%#######################################################################
\subsection{Blockkommentare in Listings}
Beachten Sie bitte, dass sich in den Listings immer mal wieder
der Orientierung und dem besseren Verständnis dienende Blockkommentare
(Kommentare über einem Block oder einer Folge von Anweisungen, wie unten
\code{// Prozess erzeugen}) finden.
In der Praxis sollte man derartige Kommentare mit Bedacht einsetzen und
 einzelne Sourcecode-Abschnitte bevorzugt in Methoden auslagern. Für die Beispiele
des Buchs dienen diese Kommentare aber als Anhaltspunkte, weil die
eingeführten oder dargestellten Sachverhalte für Sie als Leser vermutlich
noch neu und ungewohnt sind.
%
%\vfill
%\pagebreak
\lstset{mathescape=true}
\vspace{-0.85mm}
\begin{lstlisting}
// $\mbox{\bfseries Prozess erzeugen}$
final String command = "sleep 60s";  // für Windows: "cmd timeout 60"
final Process sleeper = Runtime.getRuntime().exec(command);
//    ...


// $\mbox{\bfseries Prozess über dessen Id in ProcessHandle wandeln }$
final ProcessHandle sleeperHandle = ProcessHandle.of(sleeper.pid()).
                                    orElseThrow(IllegalStateException::new);
//    ...
\end{lstlisting}
\vspace{-0.85mm}
\lstset{mathescape=false}






%#######################################################################
%#######################################################################
\subsection{Gedanken zur Formatierung}\index{Formatierung!Gedanken zur}


Die in den Listings genutzte Formatierung weicht leicht von den Coding Conventions von\index{Coding Conventions}
Oracle\footnote{\url{https://www.oracle.com/java/technologies/javase/codeconventions-introduction.html}} ab.
Ich bevorzuge diejenigen von Scott Ambler\footnote{\url{http://www.ambysoft.com/essays/javaCodingStandards.html}},
weil er insbesondere öffnende geschweifte Klammern in jeweils eigenen Zeilen vorschlägt.
Dazu habe ich ein spezielles Format namens \code{Michaelis\_CodeFormat} erstellt.
Dieses ist im Projekt-Download integriert und wird nachfolgend für eine Klasse demonstriert:


\lstset{mathescape=true}
\vspace{-0.75mm}
\begin{lstlisting}
import java.util.Arrays;
import org.apache.log4j.Logger;


public final class FormattingExample
{
    private static final Logger log = Logger.getLogger("FormattingExample");


    public static String asHex(final byte[] telegram)
    {
        log.info("asHex(" + Arrays.toString(telegram) + ")");


        final StringBuffer sb = new StringBuffer("0x");


        for (int i = 0; i < telegram.length; i++)
        {
            final String hex = Integer.toHexString(telegram[i]);
            sb.append(hex);
        }


        return sb.toString();
    }
    // ...
}
\end{lstlisting}
\vspace{-0.75mm}
\lstset{mathescape=false}








\vfill
\pagebreak
% ######################################################################
% ######################################################################
% ######################################################################
\section{Installation von Java 25 LTS}
\label{sec:Installation von Java 25 LTS}
\index{Installation}


Damit Sie die im Buch beschriebenen Java-Programme ausführen können,
müssen Sie ein JDK (Java Development Kit) installieren oder bereits installiert haben.
Beginnen wir also mit der Installation von Java.




% ######################################################################
% ######################################################################
\subsection{Java-Download}
\index{Download}\index{Java!Download}


Java ist kostenfrei auf der Oracle-Webseite \url{https://www.oracle.com/java/technologies/downloads/}
verfügbar (vgl. Abbildung \ref{Java-Download-Seite}).


% 2728?×?1858
\begin{figure}[htb]
\centering
\includegraphics[bb=0 0 1964.16 1337.76, width=\textwidth]{Abbildungen/Java25/Oracle-Java-25-Download.png}
\caption{Java-Download-Seite}\label{Java-Download-Seite}
\end{figure}


\medskip
\medskip
\noindent
Im unteren Bereich der Webseite finden Sie verschiedene Links für unterschiedliche
Betriebssysteme. Wählen Sie den für Sie passenden Link und laden Sie
die entsprechende Installationsdatei herunter, in obiger Abbildung exemplarisch für
macOS gezeigt.




% ######################################################################
% ######################################################################
\subsection{Installation des JDKs}
\index{Installation}\index{Java!Installation}


\paragraph{Windows}
Für Windows doppelklicken Sie bitte auf die \code{.exe}-Datei. Diese kann nach
erfolgreicher Installation gelöscht werden. Führen Sie das heruntergeladene Installationsprogramm aus
(z.\,B. \code{jdk-25\_windows-x64\_bin.exe}). Damit wird Java
standard\-mäßig in das Verzeichnis \code{C:\string\Program} \code{Files\string\Java\string\jdk-25}
installiert, wobei der Verzeichnisname von der gewählten Java-Version abhängt.
Akzeptieren Sie die Standard"-einstellungen und befolgen Sie die Anweisungen
während der Installation.


\paragraph{macOS}
Zum Start der Installation doppelklicken Sie unter macOS  auf die \code{.dmg}-Datei
 und folgen den Aufforderungen.  Möglicherweise müssen Sie das Administrator-Passwort
eingeben, um fortzufahren.
Nach Abschluss der Installation können Sie die \code{.dmg}-Datei löschen,
um Speicherplatz zu sparen.




% ######################################################################
% ######################################################################
\subsection{Installationsnacharbeiten}
\label{subsec:Installationsnacharbeiten}


Damit Java bei Ihnen nach dem Download und der Installation  korrekt funktioniert, sind ein paar Nacharbeiten nötig.
Dazu sollten wir die Installationsverzeichnisse sowie die der Executables zur leichteren Handhabung in
die Umgebungsvariable \code{PATH}  aufnehmen. Dies wird im Anschluss für die Betriebssysteme
Windows und macOS beschrieben. Für Unix-Derivate finden Sie weitere Informationen
hier: \url{https://www.java.com/de/download/help/download\_options.html}.
Für Windows und macOS gibt es dort ebenfalls einige ergänzende Informationen zur
Installation.




% ######################################################################
\subsubsection{Nacharbeiten für Windows}


Bei einer deutschen Variante von Windows können Sie die Umgebungsvariable \code{PATH}
unter \frqq Umgebungsvariablen\flqq{} ändern. Drücken Sie die Win-Taste
und geben Sie dann  \frqq umgeb\flqq{}  ein, bis \frqq Systemumgebungsvariablen bearbeiten\flqq{}
erscheint (vgl. Abbildung \ref{fig:Systemumgebungsvariablen}).


%800?×?750
\begin{figure}[htb]
\centering
\includegraphics[bb=0 0 576 540,height=8.5cm]{Abbildungen/Java22/Systemumgebungsvariablen.png}
\captionAndLabel{Systemumgebungsvariablen ändern}{fig:Systemumgebungsvariablen}
\end{figure}


\noindent
Mit \frqq Enter\flqq{}-Taste öffnet sich der Dialog \frqq Systemvariablen\flqq{} (vgl. Abbildung \ref{fig:Java-Standard-Verzeichnis}).

% 828?×?766
\begin{figure}[htb]
\centering
\includegraphics[bb=0 0 596.16 551.52,height=8cm]{Abbildungen/Java25/set-JAVA_HOME-Windows-jdk-25.png}
\captionAndLabel{Umgebungsvariable \texttt{JAVA\_HOME} setzen}{fig:Java-Standard-Verzeichnis}
\end{figure}


\noindent
Klicken Sie dann auf den Button \frqq Bearbeiten\flqq{} zum Öffnen
eines Bearbeitungsdialogs. Fügen Sie eine Umgebungs\-variable namens \code{JAVA\_HOME} hinzu.
Diese muss auf das Installationsverzeichnis, etwa \code{C:\string\Program} \code{Files\string\Java\string\jdk-25},
verweisen, wie es Abbildung \ref{fig:Umgebungsvariablen-1} zeigt.


% 831?×?390
\begin{figure}[htb]
\centering
\includegraphics[bb=0 0 598.32 280.8, width=\textwidth]{Abbildungen/Java25/Umgebungsvariablen-1.png}
\caption{Umgebungsvariablen bearbeiten}\label{fig:Umgebungsvariablen-1}
\end{figure}


\noindent
Ergänzen Sie die Variable \code{PATH} um den Eintrag \code{\%JAVA\_HOME\%\string\bin},
wie  es nachfolgende Abbildung \ref{fig:Umgebungsvariablen-2} zeigt.
Außerdem sollte der Eintrag möglichst weit oben stehen, damit dieser sicher den Pfad zur
aktuellen Java-Version festlegt:


% 595?×?206
\begin{figure}[htb]
\centering
\includegraphics[bb=0 0 428.4 148.32, width=\textwidth]{Abbildungen/Java21/Umgebungsvariablen-2.png}
\caption{Umgebungsvariablen ordnen}\label{fig:Umgebungsvariablen-2}
\end{figure}


\noindent
Beachten Sie bitte noch Folgendes:
Bestätigen Sie die  gesamten Dialoge immer mit OK, sodass die
Variablen gesetzt sind. Eventuell geöffnete Eingabeaufforderungen müssen
geschlossen und neu geöffnet werden, um die geänderten Variablen
wirksam werden zu lassen.


Ähnliches gilt für einige IDEs, die ebenfalls Umgebungsvariablen referenzieren und nutzen.
Beispielsweise cached IntelliJ das Environment. Bei Änderungen erfordert dies gegebenenfalls einen Neustart.






% ######################################################################
\subsubsection{Nacharbeiten für macOS}


Auch unter macOS empfiehlt es sich, einen Verweis auf Java im Pfad  in der jeweiligen
Shell (dem Terminal) passend zu setzen bzw. in das Startskript Ihrer Shell einzutragen, etwa
\code{\string~/.bash\_profile} oder neuer \code{\string~/.zshrc}:


\begin{lstlisting}
export JAVA_HOME=/Library/Java/JavaVirtualMachines/jdk-25.jdk/Contents/Home
export PATH=$JAVA_HOME/bin:$PATH
\end{lstlisting}


Für macOS geschieht dies über das Editieren einer Konfigurationsdatei wie folgt:


 \vspace{-0.75mm}
 \lstset{deletekeywords={open}}
\begin{lstlisting}
$ open ~/.zshrc
\end{lstlisting}
\vspace{-0.75mm}
\lstset{morekeywords={open}}
\lstset{mathescape=false}


Dadurch öffnet sich ein Editorfenster mit dem Inhalt der Datei. Es empfiehlt sich, dort
verschiedene Versionswechselbefehle zu definieren, um komfortabel zwischen den
Java-Versionen umschalten zu können. Für die drei Teile des Buchs sind
Java  17 LTS, Java 21 LTS und Java 25 LTS relevant.


\vfill
\pagebreak
Zum Umschalten zwischen diesen fügen Sie bitte  am Ende der Konfigurationsdatei folgende Zeilen ein:


\vspace{-0.5mm}
\begin{lstlisting}
setJdk17()
{
   export JAVA_HOME=/Library/Java/JavaVirtualMachines/jdk-17.jdk/Contents/Home
   export PATH=$JAVA_HOME/bin:$PATH
}


setJdk21()
{
   export JAVA_HOME=/Library/Java/JavaVirtualMachines/jdk-21.jdk/Contents/Home
   export PATH=$JAVA_HOME/bin:$PATH
}


setJdk25()
{
   export JAVA_HOME=/Library/Java/JavaVirtualMachines/jdk-25.jdk/Contents/Home
   export PATH=$JAVA_HOME/bin:$PATH
}
\end{lstlisting}
\vspace{-0.5mm}
\lstset{mathescape=false}


Aktivieren Sie Java 25 LTS dann wie folgt:


 \vspace{-0.75mm}
\begin{lstlisting}
$ setJdk25
\end{lstlisting}
\vspace{-0.75mm}
\lstset{mathescape=false}


Zum Zurückschalten auf beispielsweise Java 21 LTS nutzen Sie:


 \vspace{-0.75mm}
\begin{lstlisting}
$ setJdk21
\end{lstlisting}
\vspace{-0.75mm}
\lstset{mathescape=false}




\begin{praxistipp}{SDKMAN}
\noindent
\begin{center}
\begin{minipage}{\textwidth-0.6cm}
Zur Verwaltung (insbesondere Installieren und Umschalten) von verschiedenen Java-Versionen
sowie anderen Programmiersprachen und Tools gibt es das Kommandozeilentool SDKMAN.
Es bietet eine bequeme Möglichkeit, zwischen unterschiedlichen installierten Programmen zu wechseln.
Weitere Informationen finden Sie unter \url{https://sdkman.io/usage/}.
\end{minipage}
\end{center}
\end{praxistipp}




% ######################################################################
% ######################################################################
\subsection{Java-Installation prüfen}


Nach dem Ausführen der obigen Schritte sollte Java 25 LTS auf Ihrem Rechner
installiert und von der Konsole startbar sein und Sie damit bereit für die nächsten Schritte.


Öffnen Sie eine Konsole und geben Sie das Kommando \code{java -}\code{-version} ein:


\begin{lstlisting}
$ java --version
java 25 2025-09-16 LTS
Java(TM) SE Runtime Environment (build 25+37-LTS-3491)
Java HotSpot(TM) 64-Bit Server VM (build 25+37-LTS-3491, mixed mode, sharing)
\end{lstlisting}


Wenn die Ausgabe ähnlich zu dieser erfolgt, dann haben Sie Java erfolgreich installiert.


Im folgenden Text nutze ich immer \code{\$} zur Kennzeichnung von Eingaben auf der Konsole, also
dem Terminal bei macOS bzw. der Windows-Eingabeaufforderung.




% ######################################################################
\subsubsection{JShell prüfen}


Prüfen Sie der Vollständigkeit halber bitte auch noch den Aufruf sowie das Beenden des
Tools JShell, das wir für im Anschluss immer wieder für das Nachvollziehen kleiner Beispiele
sinnvoll nutzen werden:


\begin{lstlisting}
$ jshell
|  Willkommen bei JShell - Version 25
|  Geben Sie für eine Einführung Folgendes ein: /help intro


jshell> /exit
|  Auf Wiedersehen
\end{lstlisting}






% #####################################################
% #####################################################
% #####################################################
\section{Konfigurationen für Build-Tools und IDEs für Java 25 LTS}
\label{sec:Konfigurationen für Build-Tools und IDEs für Java 25 LTS}
\index{Build-Tools und IDEs!Konfigurationen für Java 25 LTS}\index{Java 25!Konfigurationen für Build-Tools und IDEs}


Build-Tools erleichtern das Kompilieren und Bündeln von Java-Klassen zu JARs. Was an Einstellungen für
Java 25 LTS zu beachten ist, beschreibe ich kurz  für die beiden wesentlichen Vertreter Gradle und Maven.
Zudem gehe ich auf Besonderheiten in der Konfiguration von IntelliJ und Eclipse ein.




% #####################################################
% #####################################################
\subsection{Java 25 LTS mit Gradle}\index{Gradle!Java 25}\index{Java 25!mit Gradle}


Zum Experimentieren mit Java 25 LTS empfiehlt sich Gradle in Version 9.1 (oder neuer).
Zudem ist es zwingend erforderlich, die Konfiguration der Java-Version nicht mehr
wie früher üblich über


\lstset{mathescape=true}
\lstset{deletekeywords={for, class}}
\vspace{-0.5mm}
\begin{lstlisting}
$\mbox{\bfseries // ACHTUNG: das ist nicht mehr korrekt! }$
$\mbox{\bfseries sourceCompatibility=25}$
$\mbox{\bfseries targetCompatibility=25}$
\end{lstlisting}
\vspace{-0.5mm}
\lstset{mathescape=false}
\lstset{morekeywords={for, class}}


zu spezifizieren, sondern wie im folgenden Listing gezeigt per Toolchain. Dazu sind ein paar, nachfolgend fett
markierte Angaben in der Datei \code{build.gradle} vorzunehmen:


\lstset{mathescape=true}
\vspace{-0.5mm}
\begin{lstlisting}
plugins {
	id 'java'
}


java {
	$\mbox{\bfseries toolchain}$ {
		$\mbox{\bfseries languageVersion = JavaLanguageVersion.of(25)}$
	}
}


// $\mbox{\bfseries Aktivierung von Preview-Features und Incubator-Modul}$
tasks.withType(JavaCompile).configureEach {
    options.compilerArgs += ["$\mbox{\bfseries -}$$\mbox{\bfseries --enable-preview}$",
                             "--add-modules", "jdk.incubator.vector"]
}
\end{lstlisting}
\vspace{-0.5mm}
\lstset{mathescape=false}


\vfill
\pagebreak
\noindent
Danach können wir wie gewohnt mit


\vspace{-0.5mm}
\begin{lstlisting}
$ gradle clean assemble
\end{lstlisting}
\vspace{-0.5mm}


ein korrespondierendes JAR bauen.


Zum Experimentieren starten Sie z.\,B. die Klasse
\code{Java25"-Runtime"-Version"-Example} aus dem Package \code{api}
aus dem gerade erstellten JAR folgendermaßen:


\begin{lstlisting}
$ java --enable-preview -cp ./build/libs/Java25Examples.jar \
       api.Java25RuntimeVersionExample
\end{lstlisting}


Wenn das Projekt wie oben mit Preview-Features kompiliert wurde, ist  es wichtig, den
JVM-Kommandozeilenparameter \code{-}\code{-enable-preview} auch beim Start anzu\-geben.
 Dann kommt es nach dem obigen erfolgreichen Gradle-Build zu folgenden Ausgaben:


\lstset{deletekeywords={for, this}}
\begin{lstlisting}
Runtime required for this: 25
latest: RELEASE_25
valueOf: RELEASE_25
25
RELEASE_25
%\end{lstlisting}
\lstset{morekeywords={for, this}}

Falls noch nicht die korrekte Java-Version verwendet wird, denken Sie vor den Aktionen daran
die Java-Version im Pfad auf Java 25 LTS umzustellen, etwa durch Aufruf des
zuvor in Abschnitt \ref{subsec:Installationsnacharbeiten}
besprochenen Kommandos \code{setJdk25}.




% #####################################################
\subsection{Java 25 LTS mit Maven}\index{Maven!Java 23}\index{Java 23!mit Maven}


Bereits mit dem im Mai 2025 aktuellen Maven 3.9.9 und auch dem im September 2025 aktuellen
 Maven 3.9.11  in Kombination mit dem Maven-Compiler-Plugin
in der Version 3.13 (oder neuer) lassen sich Java-25-Projekte ohne Probleme übersetzen.
Insgesamt setzt sich die rühmliche Tradition von Maven fort, schneller und
einfacher als Gradle neuere Java-Versionen zu unterstützen.


Um Java 25 LTS mit Maven zu nutzen und Preview-Features sowie das Incubator-Modul für
das Vector API zu aktivieren bzw. einzubinden, muss man folgende fett markierte Angaben in der Datei
\code{pom.xml} in der Sektion \code{plugins} vornehmen:


\lstset{mathescape=true}
\begin{lstlisting}
<plugins>
	<plugin>
		<artifactId>maven-compiler-plugin</artifactId>
		<version>3.13.0</version>
		<configuration>
			$\mbox{\bfseries <release>25</release>}$
			$\mbox{\bfseries <!-}$$\mbox{\bfseries- Aktivierung von Preview-Features -}$$\mbox{\bfseries ->}$
			$\mbox{\bfseries <compilerArgs>}$
                <arg>$\mbox{\bfseries -enable-preview}$</arg>
                <arg>--add-modules</arg>
                <arg>jdk.incubator.vector</arg>
            </compilerArgs>
		</configuration>
	</plugin>
</plugins>
\end{lstlisting}
\lstset{mathescape=false}


\noindent
Mit diesen Modifikationen lässt sich ein Maven-Build mit


\lstset{deletekeywords={package}}
\begin{lstlisting}
$ mvn clean package
\end{lstlisting}
\lstset{morekeywords={package}}


ausführen und es wird ein korrespondierendes JAR  -- wie von Maven gewohnt -- im
Ordner \code{target} erzeugt.


Betrachten wir als Beispiel erneut den Start der Klasse \code{Java25Runtime"-Version"-Example} aus dem Package
\code{api} -- nach dem obigen erfolgreichen Maven-Build kommt es zu den Ausgaben wie für Gradle gezeigt.


\lstset{deletekeywords={for, this}}
\begin{lstlisting}
$ java --enable-preview -cp target/Java25Examples-1.0.jar \
       api.Java25RuntimeVersionExample
\end{lstlisting}
\lstset{morekeywords={for, this}}


Das Gleiche wie für Gradle gilt auch für den Maven-Build:
Wenn das Projekt mit Preview-Features, wie in unserer obigen Build-Datei spezifiziert,
kompiliert wurde, ist  es wichtig -- und wie im obigen Aufruf geschehen --,
beim Start eines Programms aus der JAR-Datei
den JVM-Kommandozeilenparameter \code{-}\code{-enable-preview} anzugeben.








% #####################################################
\subsection{Java 25 LTS mit IntelliJ}\index{IntelliJ!Java 25}\index{Java 25!mit IntelliJ}


Um mit Java 25 LTS zu experimentieren, nutzen Sie bitte das im Sommer 2025 erschienene
IntelliJ IDEA 2025.2 oder aktueller.
Allerdings sind noch ein paar Einstellungen vorzunehmen und folgende Dinge zu beachten:


\begin{itemize}
\item Im Dialog \frqq Project Structure\flqq{}
			muss man im Bereich  \frqq \code{Project SDK}\flqq{}
			den Wert 25 auswählen und zum Experimentieren mit  Preview-Features im Bereich
			\frqq \code{Language level}\flqq{}  den Wert
			\frqq \code{25 (Preview) - Primitive Types in Patterns etc.}\flqq{} einstellen.


% 1628?×?768
\begin{figure}[htb]
\centering
\includegraphics[bb=0 0 1172.16 552.96,width=\textwidth-1mm]{Abbildungen/Java25/Java25-ProjectStructure.png}
\captionAndLabel{Einstellungen im Dialog Project Structure für Java 25 LTS}{fig:ProjectStructureJava25SettingsNormal}
\end{figure}


\vfill
\pagebreak
\item  Durch die Nutzung von Preview-Features innerhalb des Projekts ist es für das Ausführen der Beispiele
		 erforderlich, in der jeweiligen Run Configuration der
		zu startenden Klassen die VM-Option \code{-}\code{-enable-preview} anzugeben. Sofern man wie zuvor im
		Language level bereits Preview-Features ausgewählt hat, ist dieser Schritt optional.


% 1864?×?1210
\begin{figure}[htb]
\centering
\includegraphics[bb=0 0 1342.08 871.2,width=\textwidth-1mm]{Abbildungen/Java25/IntelliJ-Run-Debug-Configurations.png}
\captionAndLabel{Einstellungen in Run/Debug-Configurations}{fig:IntelliJ-Run-Debug-Configurations}
\end{figure}



\end{itemize}
\vspace{-2mm}








\vfill
\pagebreak
% #####################################################
% #####################################################
\subsection{Java 25 LTS mit Eclipse}\index{Eclipse!Java 25}\index{Java 25!mit Eclipse}
\label{subsec:Java 25 LTS mit Eclipse}


Ab Version 2025-09 wird Java 25 LTS durch Eclipse unterstützt, sofern wie nachfolgend
beschrieben das passende Plugin installiert wurde.




% #####################################################
\subsubsection{Java-25-Plugin}


In Eclipse  in der Version 2025-09 ist der Support für Java 25 LTS bereits vorbereitet und steht
durch die kurze Zeitspanne zwischen der Veröffentlichung von Java 25 LTS und Eclipse 2025-09
zunächst  als Plugin im  Eclipse Marketplace zum Nachinstallieren bereit,
wie es in Abbildung \ref{fig:EclipseMarketplaceJava25Plugin} angedeutet ist.


%1408?×?1080
\begin{figure}[htb]
\centering
\includegraphics[bb=0 0 1013.76 777.6,width=\textwidth-1mm]{Abbildungen/Java25/Eclipse-Java-25-Plugin.png}
\captionAndLabel{Installation des Java-25-Plugins aus dem Eclipse Marketplace}{fig:EclipseMarketplaceJava25Plugin}
\end{figure}


\vfill
\pagebreak
% #####################################################
\subsubsection{Compiler-Einstellungen}
Selbstverständlich müssen Sie den Compiler auf Java 25 LTS
einstellen und die Preview-Features aktivieren, um alle Beispiele ausprobieren zu können.


% 1402?×?1126
\begin{figure}[htb]
\centering
\includegraphics[bb=0 0 1009.44 810.72,width=\textwidth-1mm]{Abbildungen/Java25/Eclipse-Java-25-Compiler-Settings.png}
\captionAndLabel{Compiler-Einstellungen in Eclipse für Java 25 LTS}{fig:Java25CompilerSettings}
\end{figure}
\clearpage


\vfill
\pagebreak
\subsubsection{Run Configuration}


 Für das Ausführen der Beispiele, die Preview-Features nutzen, ist es
 erforderlich, in der Run Configuration der
zu startenden Klasse \code{-}\code{-enable-preview}  anzugeben. Das ist für
die Klasse \code{FirstVectorExample} inklusive Aktivierung des zugehörigen Incubator-Moduls
\code{-}\code{-add-modules jdk.incubator.vector} in Abbildung \ref{fig:Eclipse-Java23-Run-Debug-Configurations}
gezeigt.

% 1226?×?1002
\begin{figure}[htb]
\centering
\includegraphics[bb=0 0 882.72 721.44,width=\textwidth-1mm]{Abbildungen/Java23/Eclipse-Java23-Run-Debug-Configurations.png}
\captionAndLabel{Einstellungen in Run/Debug-Configurations}{fig:Eclipse-Java23-Run-Debug-Configurations}
\end{figure}










\clearpage
\vfill
\pagebreak
% #########################################################
% #########################################################
\section{Ausprobieren der Beispiele und Lösungen}\index{Beispiele ausprobieren}


Vielfach lassen sich die abgebildeten Sourcecode-Schnipsel in eine Klasse
mit einer \code{main()}-Methode integrieren  und danach  ausführen.
 Alter\-nativ bieten sich verschiedene Formen der Ausführung auf der Kommandozeile an.
Dazu erfahren Sie gleich mehr.
Ergänzend finden Sie alle relevanten Programme online auf
\url{https://github.com/Michaeli71/Modern-Java-25-LTS}
als Begleitprojekt zu diesem
Buch. IntelliJ erlaubt den Import als Gradle- oder Maven-Projekt.


Grundsätzlich verwende ich möglichst nachvollziehbare Konstrukte und verzichte
auf besonders ausgefallene Syntax- oder API-Features. Sofern nicht explizit im Text erwähnt, sollten
Sie die Beispiele und Lösungen daher mit der aktuellen LTS-Version Java 25 ausprobieren können.
Ab und an setze ich Preview-Features aus Java 25 LTS  ein.
Wie zuvor gezeigt, erfordert dies in den IDEs und Build-Tools gewisse Parametrierungen
sowohl beim Kompilieren als auch beim Ausführen.






% #####################################################
\subsection{Ausprobieren neuer Java-Features mit der JShell oder der Kommandozeile}
\label{subsec:Ausprobieren neuer Java-Features mit der JShell oder der Kommandozeile}


Durch die recht eng getakteten Releasezyklen von 6 Monaten ist  die Tool-Landschaft mitunter noch nicht
immer parat, wenn man schon mal mit Vorabversionen experimentieren möchte. Dann bieten sich
folgende Alternativen an:


\begin{enumerate}
\item JShell -- Die JShell ist  ein REPL-Kommandozeilen-Tool (Read-Eval-Print-Loop)\index{REPL}\index{Read-Eval-Print-Loop}
          und seit Java 9 	Bestandteil des JDKs. Damit lassen sich kleine  Sourcecode-Schnipseln schnell und einfach
          ausführen, was beim Entdecken und Ausprobieren von Neuerungen hilfreich ist.
           Eine kurze Einführung liefert  Abschnitt \ref{sec:Java + REPL}.

\item Kommandozeile -- Seit Java 11 LTS gibt es ein Feature namens Direct Compilation, mit dem sich
		einzelne Java-Dateien direkt durch Aufruf von \code{java},
		also ohne explizites vorheriges Kompilieren, starten lassen
		 (vgl. Abschnitt \ref{sec:Launch Single-File Source-Code Programs}).
		 		 Mit Java 22 und JEP 458 (vgl. Abschnitt \ref{sec:JEP 458: Launch Multi-File Source-Code
Programs}) wurde die Restriktion aufgehoben und es lassen sich nun mehrere Klassen direkt kompilieren.

\item Kommandozeile  -- Selbstverständlich ist es  immer noch möglich, mit \code{javac} und \code{java} zu arbeiten,
         wenn auch manchmal etwas mühsam -- deswegen gehe ich hier nicht näher darauf ein.
\end{enumerate}


Vielfach können Sie die abgebildeten Sourcecode-Schnipsel einfach in der JShell oder mithilfe
eines Texteditors bzw. besser noch in einer IDE eingeben. Je nach Situation empfiehlt
sich die eine oder andere Art mehr. Betrachten wir das nun genauer.






% ###############################################################
\subsubsection{Ausführen mit der \texttt{jshell}}


Mit der JShell wird es möglich, etwas Java-Sourcecode zu schreiben und Dinge schnell auszuprobieren,
ohne dafür die IDE starten und ein Projekt anlegen zu müssen.
Außerdem kann man in der JShell die neuesten Features nutzen.


Hier folgen ein paar einführende Beispiele, um Ihnen das Ausführen einzelner oder weniger
Kommandos zu verdeutlichen.
Etwas detaillierter gehe ich in Abschnitt \ref{sec:Java + REPL} auf die Möglichkeiten ein.


Starten wir die \code{jshell} und probieren einige Aktionen und Berechnungen aus.
Dabei dient eine Abwandlung eines Hello-World-Beispiels als Startpunkt:


\vspace{-0.75mm}
\begin{lstlisting}
$ jshell
|  Willkommen bei JShell - Version 25
|  Geben Sie für eine Einführung Folgendes ein: /help intro


jshell> System.out.println("Hello JShell")
Hello JShell
\end{lstlisting}
\vspace{-0.75mm}




% #######################################################################
\begin{praxistipp}{Besonderheit Preview-Features in der JShell nutzen}
\begin{center}
\begin{minipage}{\textwidth-0.6cm}
Analog zum Vorgehen beim Kompilieren muss die Möglichkeit, Preview-Featues zu nutzen,
über folgenden Angabe beim Start aktiviert werden:\index{JShell}


\vspace{-0.75mm}
\begin{lstlisting}
$ jshell --enable-preview
\end{lstlisting}
\vspace{-3mm}
\end{minipage}
\end{center}
\end{praxistipp}


\noindent
Sie können  beliebige Java-Anweisungen in der JShell eintippen.
Das ist normalerweise dadurch angedeutet, dass die einzugebenden Zeilen
mit \code{jshell>} und die nachfolgenden Zeilen bis zum Abschluss der Eingabe mit \code{...>} markiert sind:


\lstset{deletekeywords={int, return, void, open, new, import, class, private, true, this, public, final}}
\vspace{-0.75mm}
\begin{lstlisting}
jshell> int multiply(int a, int b) {
   ...>     return a * b;
   ...> }
|  created method multiply(int,int)


jshell> multiply(7, 2)
$3 ==> 14
\end{lstlisting}
\vspace{-0.75mm}
\lstset{morekeywords={int, return, void, open, new, import, class, private, true, this, public, final}}


Aufgrund der Neuerung Module Import Declarations und den Compact Source Files inklusive dem
automatischen Einbinden des Moduls \code{java.base} sind bereits diverse Funktionalitäten, wie
File-IO, Stream-API, Date and Time API usw. direkt ohne Import zugänglich.


Weitere spezielle Java-Funktionalitäten lassen sich wie gewohnt per \code{import} einbinden, hier
für die HTTP-Verarbeitung gezeigt:


\lstset{deletekeywords={int, return, void, open, new, class, private, true, this, public, final}}
\vspace{-0.75mm}
\begin{lstlisting}
jshell> import java.net.http.*


jshell> var request = HttpRequest.newBuilder()
   ...>                          .uri(new URI("https://postman-echo.com/get"))
   ...>                          .GET()
   ...>                          .build();
request ==> https://postman-echo.com/get GET


jshell> var httpClient = HttpClient.newHttpClient();
httpClient ==> jdk.internal.net.http.HttpClientImpl@78c03f1f(1)


jshell> var response = httpClient.send(request,
   ...>                                HttpResponse.BodyHandlers.ofString());
response ==> (GET https://postman-echo.com/get) 200
\end{lstlisting}
\vspace{-0.75mm}
\lstset{morekeywords={int, return, void, open, new, class, private, true, this, public, final}}


Dieses Vorgehen praktiziere ich bis maximal ca. 10 bis 20 Zeilen. Danach empfiehlt es sich,
die Zeilen mithilfe eines Texteditors einzugeben und mit der im Anschluss beschriebenen Direct Compilation auszuführen.
Noch komfortabler geht das Ganze  mithilfe einer IDE.


Manchmal möchte ich konzeptionell etwas verdeut\-lichen und
verzichte dann auf die zuvor gezeigte exakte Darstellung in der JShell, weil
die Anweisungen durch den Mix von Protokollierung mit  \code{jshell>} sowie \code{...>} und den
Zwischen\-ergebnissen unübersichtlich würden. In den Fällen greife ich auf folgende verkürzte Darstellung zurück:


\vspace{-0.75mm}
\begin{lstlisting}
int minutesPerHour = 60;
int daysPerYear = 365;


var result = Stream.of("Maria", "Mike", "Jim", "Micha").
    filter(str -> str.length() > 4).
    toList();
IO.println("filtered content: " + result);
\end{lstlisting}
\vspace{-0.75mm}


Derartige Anweisungen können nacheinander in die JShell eingegeben oder aber in einer
Java-Datei in eine \code{main()}-Methode eingebettet werden, etwa wie folgt
mit den Neuerungen aus Java 25 LTS -- dies wird auch im Anschluss nochmals aufgegriffen.


\vspace{-0.75mm}
\begin{lstlisting}
void main()
{
    int minutesPerHour = 60;
    int daysPerYear = 365;


    var result = Stream.of("Maria", "Mike", "Jim", "Micha").
           filter(str -> str.length() > 4).
           toList();
    IO.println("filtered content: " + result);
}
\end{lstlisting}
\vspace{-0.75mm}


% ###############################################################
\subsubsection{Ausführen mit der Kommandozeile -- Direct Compilation}


Werden die Programme länger, so wird das Editieren mit der JShell mühsam.
Hier kann ein Texteditor Abhilfe schaffen. Programme sind dann ohne  \code{jshell>}
abgebildet.
Als Beispiel dient ein einfaches aus einer Klasse bestehendes Programm,
das  lediglich prüft, ob zwei gegebene Zeitdauern von 20 und -12 Tagen negativ sind:


\vspace{-0.75mm}
\begin{lstlisting}
package intro;


import java.time.Duration;
import java.time.temporal.ChronoUnit;


public class DurationIntroExample
{
    public static void main(final String[] args)
    {
        System.out.println(Duration.of(20L, ChronoUnit.DAYS).isNegative());
        System.out.println(Duration.of(-12L, ChronoUnit.DAYS).isNegative());
    }
}
\end{lstlisting}
\vspace{-0.75mm}


Diese Zeilen müssen dann so eingegeben und als Datei mit dem exakt gleichen
Namen wie die Klasse inklusive der Endung \code{.java}, hier also
\code{DurationIntroExample.java},
abgespeichert werden.


Werfen wir einen kurzen Blick auf die angenehmen Neuerungen aus Java 25 LTS
und wie sich das obige Beispiel dadurch elegant, kurz und knackig schreiben lässt:


\vspace{-0.75mm}
\begin{lstlisting}
void main()
{
   IO.println(Duration.of(20L, ChronoUnit.DAYS).isNegative());
   IO.println(Duration.of(-12L, ChronoUnit.DAYS).isNegative());
}
\end{lstlisting}
\vspace{-0.75mm}


Interessanterweise benötigt es hier keine expliziten Imports mehr, da
die verwendeten Funktionalitäten aus dem Modul \code{java.base} stammen.
Dieses wird für diese kompakte Schreibweise automatisch importiert.
Details zu der Neuerung Module Imports Declarations finden
Sine in Abschnitt \ref{sec:JEP 511: Module Import Declarations}.


Nochmals zur kompakten Schreibweise. Diese nennt sich
Compact Source Files and Instance Main Methods
 (vgl. Abschnitt \ref{sec:JEP 512: Compact Source Files and Instance Main Methods}).
In dem Kontext wurde auch die Klasse \code{java.lang.IO} eingeführt, die Interaktionen
mit der Konsole leichter als früher gestaltet.


Wie bereits erwähnt, kann man (einzelne) Java-Dateien durch Aufruf von \code{java},
ohne explizites vorheriges Kompilieren, von der Konsole starten, nachfolgend unter
der Annahme einer Speicherung in einem Ordner \code{src/main/java/intro/}:


\vspace{-0.75mm}
\begin{lstlisting}
$ java src/main/java/intro/DurationIntroExample.java
\end{lstlisting}
\vspace{-0.75mm}


\noindent
Dann erhalten wir diese Ausgabe:


\lstset{deletekeywords={true, false}}
\vspace{-0.5mm}
\begin{lstlisting}
false
true
\end{lstlisting}
\vspace{-0.5mm}
\lstset{morekeywords={true, false}}








% #######################################################################
\begin{praxistipp}{Besonderheit Preview-Features}
\begin{center}
\begin{minipage}{\textwidth-0.6cm}
Bei der Direct Compilation ist für Preview-Features neben
 \code{-}\code{-enable-preview}  die Angabe einer
Source-Version mit  \code{-}\code{-source} wie folgt erforderlich --
seit Java 23 kann man darauf sogar verzichten.


\vspace{-0.75mm}
\begin{lstlisting}
$ java --enable-preview --source 21 <YOUR_CLASS_NAME>.java
\end{lstlisting}
\vspace{-3mm}
\end{minipage}
\end{center}
\end{praxistipp}
\vspace{-0.75mm}








%\vfill
%\pagebreak
% #####################################################
% #####################################################
\subsection{Ausprobieren neuer Java-Features in einer Sandbox}\index{Sandbox}


Durch die häufigen Releasewechsel und die ständig neu zu installierenden Versionen ist man
mitunter noch nicht auf dem aktuellen Stand oder vielleicht möchte man schon einmal ein
Feature ausprobieren, bevor das offizielle Release erschienen ist. Genau dafür
bietet die Seite \url{https://javaalmanac.io/} eine sogenannte Sandbox. Dort können Sie
in einem Editorfenster bereits die neuen (bis zu dem Zeitpunkt bereitgestellten) Features
ausprobieren, indem Sie kleine Programme in Form einer Klasse oder wie hier bereits als implizit deklarierte Klasse
(vgl. Abschnitt \ref{sec:JEP 512: Compact Source Files and Instance Main Methods}) und mit einer
simplen \code{main()}-Methode erstellen und dann dort ausführen, wie es Abbildung \ref{fig:Java25-Sandbox} zeigt.
Zum Experimentieren können Sie die gezeigten Zeilen als Anregung nutzen.


% 1614?×?432, => *.72
\begin{figure}[htb]
\centering
\includegraphics[bb=0 0 1162.08 311.04,width=\textwidth-1mm]{Abbildungen/Java25/Java25-Sandbox.png}
\caption{Sandbox zum Ausprobieren von neuen Funktionalitäten}\label{fig:Java25-Sandbox}
\end{figure}










% #########################################################
\section*{Entdeckungsreise zu modernem Java -- Wünsche an die Leserinnen und Leser}


Genug der Vorrede: Ich wünsche Ihnen nun viel Freude mit diesem Buch sowie einige neue
Erkenntnisse und viel Spaß beim Experimentieren mit Java 21 LTS und
 dem brandaktuellen Java 25 LTS.
Möge Ihnen der Umstieg auf die neue(n) Java-Version(en)  oder die Migration
bestehender Anwendungen  durch die Lektüre meines Buchs  leichter fallen.


Wenn Sie zunächst eine Auffrischung Ihres Wissens zu wesentlichen
Neuerungen aus Java 8 LTS bis 11 LTS benötigen, bietet sich ein Blick in den Anhang
\ref{cha:Wesentliches aus Java 8 LTS, 9, 10 und 11 LTS} an.


Teil I führt Sie kompakt in die Neuerungen aus Java 12 bis 17 LTS ein.
Danach beschreibt Teil II ausführlich die Erweiterungen aus Java 18 bis 21 LTS.
Im abschließenden Teil III werden die Neuerungen in Java 22 bis 25 LTS  im Detail behandelt.






